\documentclass[a4paper]{article}
% Español (México) con XeLaTeX: fontspec para fuentes Unicode, polyglossia
% para la división silábica y los nombres del documento en la variante mexicana.
\usepackage{fontspec}
\usepackage{polyglossia}
\setdefaultlanguage[variant=mexican]{spanish}
% Los nombres chinos van como «pinyin (original)»: xeCJK da a los caracteres
% Han su propia fuente; sin ella XeLaTeX los omite con un aviso.
% FandolSong no cubre algunos caracteres de nombres (U+73FA); WenQuanYi Zen Hei sí.
\usepackage[AutoFallBack=true]{xeCJK}
\setCJKmainfont{FandolSong-Regular.otf}
\setCJKfallbackfamilyfont{\CJKrmdefault}{WenQuanYi Zen Hei}
% polyglossia no traduce los nombres de listings.
\AtBeginDocument{\providecommand{\lstlistingname}{}\renewcommand{\lstlistingname}{Listado}%
  \providecommand{\lstlistlistingname}{}\renewcommand{\lstlistlistingname}{Índice de listados}}
\usepackage{amsmath, amssymb}
\usepackage{graphicx}
\usepackage[margin=2.5cm]{geometry}
\usepackage[most]{tcolorbox}
\usepackage{etoolbox}
\usepackage{listings}
\usepackage{hyperref}
\usepackage{booktabs}
\usepackage{float}
\newtcolorbox{knowledgebox}[1]{enhanced, colback=blue!5!white, colframe=blue!75!black, colbacktitle=blue!75!black, coltitle=white, fonttitle=\bfseries, title={#1}, attach boxed title to top left={yshift=-2mm, xshift=2mm}, boxrule=1pt, sharp corners}
\newtcolorbox{importantbox}[1]{enhanced, colback=yellow!10!white, colframe=yellow!80!black, colbacktitle=yellow!80!black, coltitle=black, fonttitle=\bfseries, title={#1}, sharp corners}
\newtcolorbox{warningbox}[1]{enhanced, colback=red!5!white, colframe=red!75!black, colbacktitle=red!75!black, coltitle=white, fonttitle=\bfseries, title={#1}, sharp corners}
\lstset{language=Python, basicstyle=\ttfamily\small, keywordstyle=\color{blue}, stringstyle=\color{red!60!black}, commentstyle=\color{green!60!black}, breaklines=true, frame=single, numbers=left, numberstyle=\tiny\color{gray}, captionpos=b, extendedchars=true}
\newcommand{\notetitle}{Explicación detallada de los parámetros de entrenamiento de veRL}
\newcommand{\noteauthors}{Elaboradas a partir de materiales públicos del curso}
\newcommand{\notedate}{2025}
\newcommand{\videochannel}{Wudaokou Nashi (五道口纳什)}
\newcommand{\videopublishdate}{2025}
\newcommand{\videoduration}{aproximadamente 25 minutos}
\newcommand{\videourl}{}
\newcommand{\videocoverpath}{cover.jpg}
\newcommand{\repourl}{https://github.com/hqhq1025/ai-course-notes}

% Footer with repo info
\usepackage{fancyhdr}
\pagestyle{fancy}
\fancyhf{}
\fancyfoot[C]{\thepage}
\fancyfoot[R]{\footnotesize\color{gray}\href{\repourl}{AI Course Notes}}
\renewcommand{\headrulewidth}{0pt}
\renewcommand{\footrulewidth}{0pt}

\begin{document}
\begin{titlepage}\centering
{\Large Notas del curso\par}\vspace{1.2cm}{\huge\bfseries \notetitle\par}\vspace{0.8cm}{\large \noteauthors\par}\vspace{0.3cm}{\large \notedate\par}\vspace{1.2cm}
\ifdefempty{\videocoverpath}{}{\includegraphics[width=0.82\textwidth,height=0.45\textheight,keepaspectratio]{\videocoverpath}\par}
\vfill
\begin{tcolorbox}[width=0.9\textwidth, colback=black!2!white, colframe=black!60, sharp corners]
\textbf{Autor o canal del video}: \videochannel\par\textbf{Fecha de publicación}: \videopublishdate\par\textbf{Duración del video}: \videoduration\par
\ifdefempty{\videourl}{}{\textbf{Enlace del video}: \href{\videourl}{\nolinkurl{\videourl}}\par}\par
\textbf{Página del proyecto}: \href{\repourl}{\nolinkurl{\repourl}}\par
\end{tcolorbox}
\end{titlepage}
\tableofcontents\newpage

\section{Introducción}
Esta entrega presenta los numerosos parámetros del script de inicio de entrenamiento de veRL, con una comprensión profunda del significado a nivel de principios de cada parámetro.

\begin{importantbox}{Importancia de comprender los parámetros}
Muchos estudiantes ya han entrenado muchos modelos, pero su comprensión del significado a nivel de principios de los parámetros es superficial. Una vez que se comprenden bien los parámetros de veRL, comprender otros marcos de trabajo como TRL también resultará muy claro.
\end{importantbox}
\section{Recursos de documentación oficial}
veRL ofrece varios documentos clave:
\begin{itemize}
  \item Documentación de Config: \texttt{ppotrainer.yaml} con toda la configuración de parámetros y su significado
  \item Documentación de Algorithm: descripción de los distintos Advantage Estimator
\end{itemize}

\section{Parámetros clave del entrenamiento}
\subsection{Parámetros relacionados con el modelo}
Incluyen configuraciones básicas como model path, tokenizer, dtype, entre otras.

\subsection{Parámetros relacionados con el algoritmo}
El clip epsilon, el KL coefficient, el learning rate, entre otros, de PPO/GRPO.

\subsection{Parámetros relacionados con el entrenamiento}
Batch size, gradient accumulation, number of rollouts, entre otros.

\subsection{Resumen de la sección}
Los parámetros se dividen en tres grandes categorías: modelo, algoritmo y entrenamiento. Entender a fondo cada parámetro ayuda a ajustar el rendimiento del entrenamiento.

\newpage
\section{Método de ajuste de parámetros: primero por capas, luego localizar el cuello de botella}
Frente a un marco de entrenamiento como veRL, el problema más común no es «demasiados pocos parámetros», sino que los parámetros no están bien separados por capas. Solo al separar el modelo, el algoritmo y la configuración de entrenamiento se puede determinar si el cuello de botella actual de desempeño proviene de la capacidad del modelo, de la estimación de la ventaja o de la orquestación del entrenamiento (lote, rollout, etc.).

\begin{knowledgebox}{Orden recomendado de ajuste de parámetros}
\begin{itemize}
  \item Primero confirmar que la configuración del modelo y del tokenizer sea correcta
  \item Después ajustar los hiperparámetros relacionados con el algoritmo, como clip, KL, learning rate
  \item Por último ajustar parámetros de sistema como lote, rollout y accumulation
\end{itemize}
\end{knowledgebox}

\subsection{Resumen de la sección}
El sentido de separar los parámetros por capas es ayudar a quien investiga a encontrar el verdadero cuello de botella con un costo de prueba y error menor, en lugar de mover todos los interruptores a la vez sin orden.

\newpage
\section{Síntesis y ampliación}
\begin{enumerate}
  \item Los parámetros de veRL se dividen en tres categorías: Config, Algorithm y Training
  \item Entender los principios de los parámetros es la base del ajuste
  \item Los principios son comunes entre frameworks; la experiencia con veRL es transferible a TRL
\end{enumerate}
\end{document}
